\section*{Overview}

Small-to-large merging is the trick I reach for when each subtree carries a container and naive merging would be
quadratic. The rule is simple: always merge the smaller container into the larger one.

\section*{When to Use It}

Use it when:

\begin{itemize}
  \item you are doing DFS on a tree,
  \item each subtree needs a set, map, frequency table, or similar container,
  \item parent states are formed by merging all child containers.
\end{itemize}

\section*{Core Idea}

After solving every child subtree, keep the largest child container as the base and insert all elements from smaller
child containers into it. That strategy limits how many times one element can move.

\section*{Key Insight}

Whenever one element moves, it moves into a container at least twice as large as before. So one element can be moved at
most \(O(\log n)\) times overall. That is the amortized reason the technique is fast.

\section*{Operations / Main Technique}

\begin{itemize}
  \item DFS the children first,
  \item select the largest child container as the main container,
  \item merge smaller child containers into it,
  \item answer the current node's query from the merged state.
\end{itemize}

\paragraph{Worked example.}
If each subtree stores the set of colors appearing in it, then small-to-large merging gives all subtree color counts in
roughly \(O(n \log n)\) instead of repeatedly rebuilding sets from scratch.

\section*{Correctness Intuition}

The merge order does not affect the final container contents, only the cost. Since all child data eventually lands in
the parent's chosen base container, the subtree summary remains correct.

\section*{Complexity Analysis}

With balanced-tree maps or sets, many common uses run in \(O(n \log^2 n)\) or \(O(n \log n)\), depending on the exact
container operations.

\section*{Implementation}

The code shows the classic subtree-color-count example using maps. The same pattern applies to sets, multisets, and
other associative containers.

\section*{Common Pitfalls}

\begin{itemize}
  \item Forgetting to swap so the larger container stays as the base.
  \item Copying containers instead of reusing pointers or references.
  \item Using the trick when the merge operation itself is too expensive or not associative enough.
\end{itemize}

\section*{Variants / Extensions}

\begin{itemize}
  \item DSU on tree, which is related but organized differently.
  \item Maintaining best frequency or color answer alongside the map.
  \item Using vectors plus coordinate compression when keys are small.
\end{itemize}

\section*{Practice Problems}

\begin{itemize}
  \item Number of distinct colors in every subtree.
  \item Frequency-based subtree queries.
  \item Merge subtree statistics where order is irrelevant.
\end{itemize}

\section*{References}

\begin{itemize}
  \item \href{https://codeforces.com/catalog}{Codeforces Catalog}.
  \item \href{https://usaco.guide/plat/merging}{USACO Guide: Small-to-Large Merging}.
\end{itemize}
